\subsection{Skill Extraction and Job-Ad Analytics}
\label{sec:sota-skill-extraction}

Online job advertisements have become a primary empirical source for observing employer demand.
Because vacancy text is unstructured, a substantial literature focuses on identifying the skills and competencies that postings require.
Khaouja~\emph{et al.} systematically surveyed 108 studies on skill identification from online job ads, documenting heterogeneity in skill bases, extraction methods, target sectors, and granularity, as well as a range of downstream goals from labor-market monitoring to curriculum design~\cite{Khaouja2021}.
Complementary work extracts soft-skill requirements and relates them to wage and demographic patterns~\cite{Calanca2019}, while domain-specific content analyses of information systems and related advertisements illustrate how technical and non-technical expectations cluster within early-career roles~\cite{Kennan2008}.

Taken together, this line of research establishes that vacancies encode actionable market signals, but that turning raw text into reliable structured fields remains a non-trivial pipeline problem.
The present work builds on that premise for the software labor market: offline enrichment is treated as a first-class engineering task whose outputs (structured metadata and comparable candidate representations) feed later statistical guidance.

\subsection{Occupation Taxonomies and Title Standardization}
\label{sec:sota-taxonomies}

A recurring obstacle in vacancy analytics is the instability of job titles: employers use overlapping labels for roles that differ in skill mix, and the same underlying occupation may appear under many surface forms.
European Skills, Competences, Qualifications and Occupations (ESCO) and related taxonomies provide standardized occupation and skill vocabularies that support comparability across postings and jurisdictions.
Recent systems such as CareerBERT combine ESCO with large vacancy corpora (e.g., EURES) to embed resumes and generic occupations in a shared space for matching~\cite{Rosenberger2025}.
Such approaches underline the value of title~/~occupation normalization as infrastructure for any downstream comparison between people and market demand.

In this thesis, standardized role information is treated as a critical metadata field for cohort-level summaries (for example, rankings of frequent normalized titles).
Whether normalization follows a fixed taxonomy, emergent clustering, or a hybrid scheme is a Method decision; the state of the art makes clear that some form of title standardization is a prerequisite for interpretable aggregates.

\subsection{Job Recommendation and Person--Job Matching}
\label{sec:sota-recommenders}

A second major stream concerns e-recruiting recommenders that retrieve or rank vacancies for candidates (and, symmetrically, candidates for recruiters).
Surveys of job recommender systems describe content-based, collaborative, and hybrid matching as responses to information overload on online recruiting platforms~\cite{AlOtaibi2012}.
Industrial and research systems increasingly rely on embedding-based retrieval and multi-stage ranking: for example, fused embeddings with approximate nearest-neighbor search for large-scale job--candidate matching~\cite{Zhao2021}, and bidirectional NLP recommenders that score resumes against postings and vice versa~\cite{Alsaif2022}.
CareerBERT similarly optimizes recommendation quality for counselors and seekers, albeit against standardized ESCO occupations rather than only individual ads~\cite{Rosenberger2025}.

These systems are successful at their stated objective---surfacing relevant \emph{items}---but that objective differs from career \emph{strategy} support.
Their typical outputs are ranked lists of jobs or occupation labels; evaluation centers on retrieval and ranking metrics (e.g., click-through or ranking quality), not on inspectable distributional summaries of market requirements conditioned on a user profile.

\subsection{From Matching to Market Intelligence: Gap and Positioning}
\label{sec:sota-gap}

Across the strands above, two capabilities are comparatively mature: (i)~extracting structured signals from vacancy text, and (ii)~matching people to jobs or taxonomy nodes via similarity search.
What remains less developed as an integrated user-facing paradigm is \emph{profile-conditioned market intelligence}: using proximity in a representation space only to select a relevant cohort of enriched offers, then returning deterministic aggregates over that cohort's metadata---frequent standardized roles, stack-fit rates, experience distributions, and related indicators---so that users can interpret the market and decide for themselves.

Conversational retrieval-augmented generation (RAG) over vacancy corpora is a plausible alternative interface, but it targets narrative question answering rather than reproducible statistical summaries; production-grade conversational guidance on top of an enriched dataset is treated here as future work rather than as the core contribution~\cite{Gao2023}.
The present thesis therefore positions itself between vacancy analytics and job recommendation: it reuses enrichment and embedding ideas from both, while defining the guidance layer as an information system of inspectable aggregates rather than as a recommender or a chatbot.
